\documentclass[11pt,reqno]{amsart}
\usepackage{amssymb,booktabs}
\usepackage[colorlinks=true,linkcolor=blue,citecolor=blue,urlcolor=blue]{hyperref}
\hypersetup{pdftitle={Majorization in odd-digit Kaprekar dynamics and the nine-digit arrival law},pdfauthor={Rustem Galyavin}}
\usepackage[margin=1.15in]{geometry}
\setlength{\emergencystretch}{1.3em}
\newtheorem{theorem}{Theorem}
\newtheorem{lemma}{Lemma}
\newcommand{\sort}{\operatorname{sort}}
\newcommand{\ii}{\operatorname{int}}
\newcommand{\law}{\mathcal L}
\newcommand{\RR}{\mathbb R}
\newcommand{\ZZ}{\mathbb Z}
\newcommand{\runcheck}[1]{\texttt{./reproduce.sh\ #1}}
\newcommand{\file}[1]{\nolinkurl{#1}}
\title[Majorization and nine-digit Kaprekar dynamics]{Majorization in odd-digit Kaprekar dynamics\protect\\ and the nine-digit arrival law}
\author{Rustem Galyavin}
\thanks{Contact: \href{mailto:rustem@galyavin.com}{rustem@galyavin.com}.
ORCID \href{https://orcid.org/0009-0006-6875-9170}{0009-0006-6875-9170}.}
\date{18 November 2026}
\subjclass[2020]{Primary 11A63; Secondary 37E99, 26D15, 60F05}
\keywords{Kaprekar map, majorization, digit differences, first arrival time, lattice index}
\begin{document}
\begin{abstract}
Arrange a string of digits in decreasing and increasing order, subtract, and repeat.
For any odd number of digits the continuous large-base map decreases the vector
of successive gaps in majorization order. Every real orbit approaches a periodic
orbit. For nine digits and strictly positive gaps, periodicity is reached in finite
time, with period one or three. For a uniformly chosen nonconstant string, the
initial differences, the time to the first cyclic number and the differences of
that number, all divided by the base, have five distinct joint limits, indexed by
the base modulo five. A uniform linear bound on the arrival time gives convergence
of all fixed positive moments and exponential moments. The nine-digit proofs use
computer-assisted affine inequalities and exact lattice counts.
\end{abstract}
\maketitle
\section{Introduction}
How much can the large-base limit of digit subtraction tell us about the time
to reach a cycle, or about the first number one reaches on that cycle?
Fix $n=2k+1$ slots, $k\ge1$, and retain leading zeros. At base $B\ge2$
let $K_B$ subtract the ascending arrangement of a string from the descending
one. For sorted digits $d_1\ge\cdots\ge d_n$, put
$q_i=d_i-d_{n+1-i}$. The resulting number is
$\sum_{i=1}^k q_i(B^{n-i}-B^{i-1})$, so there is an induced map $F_B$
on integer difference vectors $q$. For an ordered vector $x$ in $[0,L]^k$ put
\[
 g_L(x)=(L-x_1,x_1-x_2,\ldots,x_{k-1}-x_k,x_k).
\]
Write $S_j(g)$ for the sum of its $j$ largest entries. For equal-sum vectors
$h\prec g$ means $S_j(h)\le S_j(g)$ at every rank.

The continuous leading map $f$ sorts the list $1,x_1,\ldots,x_k,
1-x_1,\ldots,1-x_k$ and takes its symmetric differences. This is the
normalized limit of the positive-inner subtraction formula below, extended
continuously to the closed simplex.
\pagebreak[3]
\begin{theorem}\label{th:odd}
For every ordered $x\in[0,1]^k$, $g_1(fx)\prec g_1(x)$. Every orbit of $f$ approaches
a periodic orbit of period at most $(k+1)!$. For integer inputs with $q_k>0$,
\[
 S_j(g_B(F_Bq))\le S_j(g_B(q))+2\qquad(1\le j\le k).
\]
\end{theorem}

Choose the initial string $N$ uniformly among the $B^n-B$ nonconstant
ordered strings. Let $T_B=\min\{t\ge0:K_B^tN\text{ is cyclic}\}$,
$X_B=q(N)/B$, and $Z_B=q(K_B^{T_B}N)/B$. The last variable retains the
phase at which the original number first enters its cycle.
\begin{theorem}\label{th:nine}
In nine digits every real leading orbit whose initial gaps are all positive
becomes exactly periodic after finitely many steps, with period one or three.
For Lebesgue-almost every initial point, the leading orbit, written in the
radius coordinates $r=\sort(|2x_i-1|)$, enters the interior of the core
polytope $\mathcal C$ of Section~\ref{sec:entry}.
For each $\beta\in\{0,1,2,3,4\}$, as $B\to\infty$ through $B\equiv\beta\pmod5$,
\[
 (X_B,T_B/B,Z_B)\ \Longrightarrow\ \mu_\beta .
\]
The five probability measures, given by \eqref{eq:law} with rational coefficients determined by the finite calculation of Section~\ref{sec:law} and listed in full in the ancillary material, are pairwise distinct.
Their terminal coordinates lie on the segment
$Q(s)=(s,2s-1,3s-2,1-s)$, $4/5\le s\le5/6$.
\end{theorem}
\begin{theorem}\label{th:clock}
In nine digits $\max_N T_B=\Theta(B)$. Along each of the five base classes,
for every bounded continuous $b(x,z)$,
\[\mathbb E[b(X_B,Z_B)\varphi(T_B/B)]\longrightarrow \int b(x,z)\varphi(t)\,d\mu_\beta\]
for
$\varphi(t)=t^p$, $p>0$, or $\varphi(t)=e^{\theta t}$, $\theta\in\RR$.
\end{theorem}
The upper bound behind $\max_N T_B=O(B)$ is explicit
(Section~\ref{sec:capture}), but its constant, of order $10^9$, is far
from sharp.

A base translation, which Thakur proposes for sufficiently large cyclic states \cite[Remark~4.3]{thakur}, explains which finite periods can occur. Put $m=k+1$,
$v=(k,k-1,\ldots,1)$. In the paired region, defined in
Section~\ref{sec:cov}, the borrowed digits occur in $k$ successive
two-digit levels below the largest digit.
\begin{theorem}\label{th:cov}
For $q$ in the paired region and every integer $j\ge0$,
$F_{B+jm}(q+jv)=F_B(q)+jv$, and the translated source is again paired.
Condition the original population on $q/B$ in a fixed Jordan-measurable
subset of the strict simplex of positive limiting weight. Suppose all but
vanishing mass have their entire difference orbit in the paired region.
If $(X_B,T_B/B,Z_B)$, so conditioned, converges in law along every residue
class of $B$ modulo $M$, then the limits depend only on the base modulo
$\gcd(M,k+1)$.
\end{theorem}

The geometry behind Theorem~\ref{th:odd} can be seen on a circle. Place
the differences and their reflections around it. Recolor the combined
points alternately. Long empty arcs become scarcer for either color, and
the new gaps are exactly the output gaps. This orders all the partial
sums at once. In nine digits it confines almost every leading orbit to
a family of three-cycles and fixed points. An actual subtraction changes two endpoint digits by one. A
three-step return that closes in the real map can therefore translate
the integer state a few units. Repetition over a distance proportional
to the base accounts for the arrival time.

When such a translation reaches a face of the periodic region, its last
step leaves a remainder. Two inputs that agree to leading order can have
different remainders and hence different speeds or different destinations.
Their proportions come from the lattice of all inputs that reached the face.
A change of base by five preserves the ordering of the levels of the nine borrowed
digits after a matching translation of the four differences. Congruences at
individual faces can have larger moduli, but this common translation survives
them all. The separation calculation in Section~\ref{sec:law} shows that the
five remaining classes really occur. Theorem~\ref{th:cov} thus explains the
modulus; it is not a further limit theorem. Thakur also exhibits nine-digit cycle
families in all five base classes \cite[Section~4.4]{thakur}.
What is new here is the paired condition, an explicit sufficient condition for
exact covariance along an orbit, and the resulting gcd restriction on
distributional limits. How far arbitrary cycles satisfy the paired condition
is left open.

Section~2 proves Theorem~\ref{th:odd}. Section~\ref{sec:entry} proves finite
entry in nine digits, Section~\ref{sec:law} derives the five laws from face
exits and lattice remainders and exhibits, in Section~4.4, one route on which
they differ explicitly, Section~\ref{sec:capture} proves the uniform
bound, and Section~\ref{sec:cov} proves Theorem~\ref{th:cov}. The appendices
state every computer-assisted step.

The difference coordinates and the borrowed digit formula appear in
Thakur \cite[Section 2.7]{thakur}. Kay and Downes-Ward studied fixed points
and cycles in odd bases \cite{kay}, and Dahl studied convergence distances,
basin statistics and projected drift for three through six slots
\cite{Dahl}. The majorization argument for a finite limit set is also used
for doubly stochastic maps by Abdulghafor et al.\ \cite[Theorem 3.3 and
Remark 3.4]{major}. The four-digit cycles and arrival times in even bases
are treated in \cite{G1,G2}. The five-digit cycle weights in \cite{G3} have
period three,
in keeping with $k+1=3$.

\section{Reflection on a circle}
For $0\le x_k\le\cdots\le x_1\le C\le L$, $L>0$, let $D_C(x)$ take
symmetric differences after sorting $C,x_i,C-x_i$. We claim
\begin{equation}\label{eq:circle}
 g_L(D_Cx)\prec g_L(x).
\end{equation}
On the circle of length $L$ consider $A=\{0,x_1,\ldots,x_k\}$ and its
reflection $C-A$, with multiplicities. Both have gaps $g_L(x)$, up to order.
Write their union as $0=p_0\le\cdots\le p_{2k+1}=C$, so
$p_i+p_{2k+1-i}=C$. Dropping $p_0$ and folding the sorted list gives
$y_j=p_{2k+2-j}-p_j$. Thus $L-y_1=L-p_{2k}$, and the remaining
gaps equal the two-step spacings $p_{i+2}-p_i$ with $0\le i\le k-1$
in some order. Indeed $y_j=C-p_{j-1}-p_j$, whence
$y_j-y_{j+1}=p_{j+1}-p_{j-1}$ and $y_k=p_{k+1}-p_{k-1}$.
Reflection identifies the spacing at $i$ with that at $2k-1-i$,
so each of these $k$ spacings occurs once among the even initial indices.
The colors of the alternate points therefore both have the gap multiset
$g_L(y)$, including their wrapping gaps. Coincidences follow by continuity.

A forward arc of length $t$ misses a set of points with gaps $h$ on
a set of starting positions of length $\sum_i(h_i-t)_+$.
An arc containing no combined point misses both colors under either coloring.
One containing exactly one misses one color. Two consecutive points already
give both alternate colors. Integration yields the inequality between these
sums for $h=g_L(y)$ and $h=g_L(x)$, for every $t$. Minimizing
$jt+\sum_i(h_i-t)_+$ gives $S_j(h)$ and proves \eqref{eq:circle}.

\begin{proof}[Proof of Theorem~\ref{th:odd}]
Take $C=L=1$. The ordered partial sums decrease to limits, so the decreasing
rearrangement of the gaps converges. Let $\omega(x)$ be the set of accumulation
points of the gap orbit. Each is a permutation of a single vector, so
$\omega(x)$ is finite, and the continuous gap map sends $\omega(x)$ onto itself.
Take disjoint small neighborhoods of its points. Eventually the orbit stays
in their union, with successive nearest-point labels following the induced
permutation by continuity. The error from these labels tends to zero.

For integer $q_k\ge1$ the subtraction digits are
\begin{equation}\label{eq:borrow}
 q_1,\ldots,q_{k-1},q_k-1,\ B-1,\ B-1-q_k,\ldots,B-1-q_2,B-q_1 .
\end{equation}
Here the reflected list before the final $B-q_1$ consists of indices
$k,k-1,\ldots,2$ and is empty for $k=1$. Relative to the list for $C=B-1$
two entries change by one. Sorting minimizes the sum of matching distances,
and each digit enters at most one symmetric difference, so the difference
vectors change by at most two in $\ell^1$. For a change $w$ of difference
vector the gap change has sum zero and norm at most $2\|w\|_1$. The sum
over any subset of gaps changes by at most $\|w\|_1$. Maximize over subsets
and apply \eqref{eq:circle} with $L=B$.
\end{proof}

\paragraph{Example.}
In base $10$ with five digits take $q=(8,3)$, realized by the string
$86530$. Its subtraction is $86530-03568=82962$.
The reflection with span $9$ gives the list $9,8,3,6,1$ and differences
$(8,5)$. The two digit changes replace $3,1$ by $2,2$ and give actual
differences $(7,6)$. The respective ambient gap vectors of the source,
reflection and actual image are $(2,5,3),(2,3,5),(3,1,6)$. The largest
gap rises by one under borrowing.

\section{Finite entry in nine digits}\label{sec:entry}
Write $r=(a,b,c,d)=\sort(|2x_i-1|)$ in decreasing order. The first image is
$C(r)=\tfrac12(1+a,a+b,b+c,c+d)$. Indeed the four larger members of the
complementary pairs are $(1+r_i)/2$, in that order. Thereafter $fC=CH$, where
\begin{align*}
 H(r)&=\sort(a,|a+b-1|,|b+c-1|,|c+d-1|),\\
 D(r)&=2g_1(Cr)=(1-a,1-b,a-c,b-d,c+d).
\end{align*}
The core is the region on which the leading map becomes a rotation of
order three. Precisely, the core $\mathcal C$ is the polytope $A(r,1)^t\ge0$ with rows
\begin{equation}\label{eq:core}
\begin{pmatrix}
-2&0&2&1&1\\-1&-1&-1&-1&2\\-1&-1&1&1&1\\0&-1&-1&0&1\\
0&0&0&1&0\\0&0&1&-1&0\\0&1&-1&0&0\\1&-1&0&0&0\\
1&0&1&1&-1\\1&2&1&0&-2\\2&0&-1&0&-1\\2&1&-2&-2&-1
\end{pmatrix}.
\end{equation}
On $C(\mathcal C)$ the update is $R(u,v',w,z)=(u,u-z,v'-z,v'-w)$,
with $R^3=I$. Put $W=\max(0,-\min_i A_i(r,1)^t)$.
Write $\mathcal E$ for the following three-cycle family, expressed in
its three successive phases:
\begin{align*}
 &(a,2a-1,2a-1,d),\quad (a,2-2a-d,3a-2,3-4a),\\
 &(a,a,1-a+d,1-a-d),\qquad 5/7\le a\le3/4,\quad 0<d\le1-a .
\end{align*}
The four affine implications of Appendix~\ref{app:entry} are the only
computer-assisted local assertions needed to prove finite entry.

By Theorem~\ref{th:odd} and invertibility of $C$, each radius orbit approaches
a periodic orbit. Its first radius is nondecreasing since $(Hr)_1=\max(a,1-c-d)$,
and the smallest gap is nondecreasing by majorization. On the limiting cycle
the sorted gaps and the first radius stay constant. Implication (i) leaves the
core, $\mathcal E$, and two exceptional orbits having a zero gap. An initially
positive minimum gap excludes the last two orbits and bounds $a$ away from one.

For a limit in $\mathcal E$, take every third point with the phase of the first
displayed vector. Write $p=b-c$, $s=c-2a+1$, $\nu=p+|s|$. The sections
eventually satisfy the inequalities of (ii), for one fixed sufficiently small
$\eta>0$, since $d$ tends to a positive value and $a$ increases to its limit.
Thus $a$ is eventually constant, including at intervening steps. Implication
(iii) now makes $\nu$ nondecreasing on late sections. Its limit is zero, so
eventually $b=c=2a-1$. Only $d$ can still vary. With $a$ fixed, its third
iterate on this plane is a continuous, finitely piecewise-affine function
of $d$ with integer slopes. Near its limit $d_*$ each applicable piece
satisfies $d'-d_*=l(d-d_*)$ with $l\in\ZZ$. A nonzero $l$ cannot shorten
the distance, so convergence forces an exact hit in finite time.

For a limit in the core, $W$ tends to zero. If the orbit never entered,
(iv) would give an arithmetic progression on every sufficiently late third
section, since the minimum of $D$ is bounded below. Boundedness forces
its displacement to vanish, contradicting convergence of the positive $W$
to zero. Thus entry is finite. Every branch of $H$ is unimodular: before
sorting it has the first coordinate row and three signed adjacent-sum rows.
The exceptional plane has codimension two, and the core boundary and sorting
ties lie on finitely many hyperplanes. Their finite preimages are null.
The same holds after pulling back under the signed folds of the initial $x$.
This proves the geometric assertions of Theorem~\ref{th:nine}.

\section{Arithmetic on the periodic family}\label{sec:law}
Let $\Omega=\{1>x_1>x_2>x_3>x_4>0\}$. For $x\in\Omega$, take the first
$r_*=H^e\sort(|2x_i-1|)$ in $\ii\mathcal C$ and put $y=C(r_*)$. The
previous section proves that $e$ is finite almost everywhere. A finite
polyhedral partition of the core, constructed below, assigns $y$ a route
$j$. On route $j$ let $(p_{j\beta l},t_{jl},U_{jl})$ denote the probabilities,
clock rows and terminal matrices obtained by the integer calculation below.
Here rows and matrices act on $\bar y=(y,1)^t$. The joint law is
\begin{equation}\label{eq:law}
 \mu_\beta=\int_\Omega 9!\prod_{i=1}^{5}g_1(x)_i
 \sum_l p_{j\beta l}\,\delta_{(x,t_{jl}\bar y,U_{jl}\bar y)}\,dx .
\end{equation}
In this formula the five gap entries are numbered from one.
The supplement gives all coefficients as rational numbers.

\subsection{Motion up to the first cycle}
In the core write $y=(u,v',w,z)$. From the interior one can reach six faces, with slacks
\[
\begin{array}{ll}
 A_0=u+z-1,&B_0=v'+w-1,\\
 A_1=u+v'-w-1,&B_1=u+v'-2z-1,\\
 A_2=2u-v'-1,&B_2=2u-2v'+2w-z-1 .
\end{array}
\]
For integer states replace the constant $1$ by $B$. In the bulk the
three-step return translates by $(-3,-2,-1,-1)$, so the $A$ slacks
decrease by four and the $B$ slacks by three. For any translation
$\delta$ with word length $\ell$, the first face hit minimizes $h/s$ over
all decreasing slacks $h$, with $s=-h_{\rm lin}\delta$. The limiting move
is $y\mapsto y+(h/s)\delta$ and its clock is $\ell h/s$. A bounded number
of ordinary steps rotates the exit to a canonical face. At the first
exit $A_i$ apply $R^{i+1\bmod3}$ and at $B_i$ apply $R^i$.

\begin{table}[ht]
\centering\small
\caption{Returns on the canonical faces. The parenthesized exponent is the
number of rotations $R$ before the next return.}\label{tab:routes}
\begin{tabular}{@{}lll@{}}
\toprule
Face & Return $(\ell,\delta)$ & Exit $\longrightarrow$ next face\\
\midrule
$A_2$ & $12,(-6,-12,-6,-4)$ & $A_0(1),A_1(0)\to AA$; $B_1(0)\to AB$\\
$B_0$ & $6,(-6,-1,1,-2)$ & $A_0(1)\to A_2$; $A_1(2)\to AB$\\
      & & $B_1(0),B_2(2)\to BB$\\
$BB$ & $12,(-12,-2,2,-7)$ & $A_1(2)\to AB$; $B_2(0)\to BBB$\\
$AA$ & $42,(-10,-20,-30,-20)$ & $A_0\to AAA$\\
     & $9,(-2,-4,-6,-3)$ & $A_0\to AAA$\\
$AB$ & $24,(-12,-24,-12,-18)$ & $A_1\to AAB$\\
\bottomrule
\end{tabular}
\end{table}
Here $AA=A_1\cap A_2$, $AB=A_2\cap B_1$, $BB=B_0\cap B_1$, with each
letter denoting the zero set of the indicated slack. The last letter of
$AAA,AAB,BBB$ indicates the additional active face, as specified by the
incoming exit. At $AAB$ the integer remainders select a release to the
$AA$ return or a nine-step drift to its endpoint. At $AAA$ they select a
stationary cycle on $Q$ or a moving return with time per unit longitudinal
decrease $16$ or $37/3$. At $BBB$ the corresponding rate is one.
These decisions and the first-capture inequalities are finite affine
calculations specified in Appendix~\ref{app:routes}.
Pulling back the strict first-exit comparisons gives 33 nonempty route
domains: twelve end at $AA$, fifteen at $AB$, and six at $BBB$ before the
terminal computations. For almost every source on these routes the
conditional joint law has respectively five, four and one distinct atoms; on an
exceptional hyperplane two of them can coincide.

\subsection{Counting incoming remainders}
Suppose the incoming states at base $B=5b+\beta$ form the affine lattice
$q=q_0(b)+M\ZZ^d$. For a return as above, reconstructing the sorting
inequalities gives a threshold $m_0$: after $N$ completed returns, at the first state where the return word fails, the slack
$h(q+N\delta,B)=e'$ lies in $[m_0-s,m_0)\cap\ZZ$. Thus
\[
 (h_{\rm lin}M)z-sN=e'-h(q_0(b),B).
\]
If $g$ is the gcd of the entries on the left, including $s$, this equation
is soluble exactly when $g$ divides its right-hand side. The residue classes
of $h_{\rm lin}Mz\pmod s$ have image of size $s/g$. Every accessible
$e'$ has conditional weight $g/s$. Unimodular column operations give a
particular solution and an integer basis for the kernel. Applying
$[M\mid\delta]$ gives the outgoing lattice with its full index.
For example, an exit from $AB$ ordinarily reaches the direction
$d_0=(2,4,6,3)$ with primitive step. After the release from $BB$ the
incoming step is $d_1=(4,8,7,6)$. Since $(A_1)_{\rm lin}d_1=5$
and $(A_1)_{\rm lin}\delta_{AB}=-24$, the outgoing step is
$24d_1+5\delta_{AB}=18d_0$. This changes the subsequent count.

On a fixed strict route each earlier remainder choice restricts the original
integer inputs to a coset of finite index. The map to the current lattice
is onto by construction. Preimages of equal-index classes therefore have
equal index, even when the current lattice has lower rank. Weighted lattice
sums on compact polytopes give equal limiting masses of these classes
locally at each macroscopic input. This proves the product rule for
conditional weights. Checking divisibility of the coefficient of $b$ at
each split proves that $B\bmod5$ suffices on every route.

\subsection{From core sources to original numbers}
The ordered-input multiplicity has leading term
$B^{k+1}\rho_n(q/B)$, where $\rho_n(x)=n!\prod_i g_1(x)_i$.
Indeed the first gap counts translations of the sorted string. Each
other gap is split into two adjacent digit spacings. When both are positive
an integer gap $d$ permits $d-1$ splits, and each resulting sorted string
has $n!$ orders. Choosing a zero spacing costs a power of $B$.
This gives an $O_n(B^k)$ remainder on the strict simplex and the density
in \eqref{eq:law}. Boundary strings have vanishing weight; equivalently
the gaps are sums of pairs of adjacent uniform order-statistic spacings.

On any compact subset of a cell prescribing a strict finite entry word and
subsequent strict route, all the prefix sorting margins are positive.
For sufficiently large $B$ borrowing changes only the affine constants
in that prefix. At each step the linear part has determinant $\pm1$:
the four signed factors of two in the radius coordinates cancel the four
halves in $C$, whose numerator matrix has determinant one. Hence the
incoming lattice for the core is $\ZZ^4$. The prefix costs bounded time and
the ensuing core trajectory has a transient state, which precludes capture
earlier in the prefix.

Let $h$ be the first cyclic difference time. Always $h\le T_B\le h+1$,
because the number following a cyclic difference is cyclic. The number
map from differences is injective: if $r$ is the final positive index,
the last digit below $B-1$ among the leading $k$ output slots is $q_r-1$
at position $r$, and the preceding digits recover $q_1,\ldots,q_{r-1}$.
For a transient difference input this gives $T_B=h+1$. Thus at finite base
the requested terminal vector is the successor of the first cyclic difference.
Every computed first cyclic difference divided by $B$ tends to some $Q(s)$.
On this compact segment the inner difference is positive, $f(Q(s))=Q(s)$
by direct sorting, and the normalized one-step borrowing error is $O(B^{-1})$.
The two phase conventions therefore have the same limiting terminal
coordinate. Compact exhaustion of the countably many strict entry cells
proves \eqref{eq:law}; their complement has Lebesgue measure zero.
No uniform bound on the entry word length is required.

\subsection{Five different measures}
For the route $A_0,B_1,AB$, put
\[
 L=7u/4-v'-z/4-3/4,\qquad s=-u/4+v'/2-z/4+3/4.
\]
Here $4/5<s<5/6$. The four elapsed times are $L+c(s-4/5)$ with the
following coefficients and probabilities. Only the first column ends at
$Q((s+4)/6)$, the remaining three end at $Q(4/5)$.
\[
\begin{array}{c|rrrr}
\beta\ \backslash\ c&7/2&9/2&50/9&37/6\\\hline
0&11/48&5/48&137/288&55/288\\
1&29/128&5/48&47/96&23/128\\
2&29/144&5/48&607/1152&193/1152\\
3&107/576&5/48&69/128&197/1152\\
4&233/1152&5/48&73/144&215/1152
\end{array}
\]
The radius $(18,14,6,3)/25$ is strictly inside the route.
A continuous test function can localize $x$ in one small inverse image
there and select the moving terminal component, which stays a positive
distance from $Q(4/5)$. The five integrals are distinct since the first
column has five distinct entries and the input density is positive. The
five base classes also differ in the probability of the third clock value; for
instance, its weights for $\beta=2$ and $\beta=4$ differ by $23/1152$. This
completes the proof of Theorem~\ref{th:nine}. A compact subcell here has a
clock bounded below by a positive constant for every class, giving the
lower bound in Theorem~\ref{th:clock}.

\section{Uniform time to capture}\label{sec:capture}
Write $g=g_B(q)$, $m_g=\min g_i$, $M_g=\max g_i$, and set
\[
 G(q)=\sum_{i<j}|g_i-g_j|,\qquad P(q)=B-4m_g-\min(M_g,2m_g).
\]
At a positive source in nine digits,
$S_i(g_B(F_Bq))\le S_i(g)+1$ for $1\le i\le4$
(Appendix~\ref{app:clock}). Thus $0\le G\le4B$ and $G(F_Bq)\le G(q)+8$,
using $G=2\sum_{i=1}^4 S_i-4B$. The descent assertion is
\begin{equation}\label{eq:descent}
 \min_{1\le i\le9}G(F_B^iq)<G(q)
\end{equation}
if $B\ge40$, $P(q)>100$, and all ten states have positive inner
difference. The step at which the decrease occurs may vary with $q$.

Assume from now on that $B\ge40$. A first visit from positive inner
difference to zero inner difference occurs within three steps, if it occurs at
all. From a state with zero inner difference its first positive-time return to
zero inner difference therefore takes at most four steps. The cover in
Appendix~\ref{app:clock} has at most $308B$ returning zero states. Until
capture none can recur. Allowing one final visit, by time $1232B+7$
the orbit is cyclic or permanently positive.

If $P\le100$, choose one maximal gap, write the other four as $m_g+e_i$,
and observe that $\sum e_i\le100$. Their positions, offsets and $m_g$
determine all five gaps. Hence there are at most $M=5(B+1)\binom{104}{4}$
such states. In the permanently positive orbit, take a single step at
each of these exceptional states. Elsewhere choose a decreasing block
of at most nine steps by \eqref{eq:descent}. Until capture there are at
most $M$ exceptional starts, and telescoping bounds the completed
decreasing blocks by $4B+8M$. A final partial block costs at most nine
steps. Adding preparation and one step for the original number gives
\[
 T_B\le1268B+17+365(B+1)\binom{104}{4}\qquad(B\ge40).
\]
Together with the lower bound above this proves the order of the maximum.
The normalized clocks now lie in one fixed compact interval. Each stated
moment follows from the joint weak limit by bounded continuous extension,
proving Theorem~\ref{th:clock}.

\section{Covariance under a change of base}\label{sec:cov}
For $q_k>0$ set $a_i=q_i-\boldsymbol1_{i=k}$ and
$b_i=B-q_{k+1-i}-\boldsymbol1_{i<k}$. The paired region is defined by
$\min(a_i,b_i)\ge\max(a_{i+1},b_{i+1})$ for $i<k$.
\begin{proof}[Proof of Theorem~\ref{th:cov}]
(Compare \cite[Remark~4.3]{thakur}.) Both entries at level $i$ increase by $jv_i$ under the stated translation,
and each adjacent level separation increases by $j$. Starting with $B-1$,
the increments of the descending digits are $j(m,k,k,k-1,k-1,\ldots,1,1)$.
Their symmetric differences give $jv$. Induction along an entirely paired
orbit translates every state, preserves the first cyclic difference time,
and preserves its successor. Initially cyclic differences have the same
translated phase convention except for inputs already cyclic as numbers;
these number at most $O(B^k)$, negligible in the original population.

Match $q$ with $q+v$ at base $B+m$. Normalized initial and terminal vectors
differ by $O(B^{-1})$. Jordan measurability and the local multiplicity
formula imply that the normalized weights of unmatched or excess inputs
tend to zero, by Riemann sums; the exceptional unpaired mass also tends
to zero. The two normalized difference clocks differ by $mh/[B(B+m)]$.
Proper limits on finitely many classes make $h/B$ tight, so this difference
vanishes in probability. The number clocks differ from the difference clocks
by at most one. Thus the limiting laws agree under a shift of $m$ modulo
$M$, proving the gcd assertion.
\end{proof}
On the nine-digit core every sorting word used in Table~\ref{tab:routes},
including its boundary and terminal words, respects these paired levels.
The direction $v/5=Q(4/5)$ lies on all six decreasing faces. The covariance
therefore explains the period allowed on the core; finite unimodular
prefixes carry the same five weights back to the other original sources.

\appendix
\section{Universal real entry implications}\label{app:entry}
For each implication start with real $1>a\ge b\ge c\ge d\ge0$,
write $r_i=H^ir$, $a_i=(r_i)_1$, and impose the defining equations for
each iterate. Encode an absolute value by its two linear
branches and implement descending sort by successive min/max comparisons.
The statements below can be tested by adjoining the negated conclusion.
\begin{enumerate}
\item[(i)] If $a_0=\cdots=a_4$, $a_4<1$, and the decreasing sorts of
$D(r_i)$ agree for $0\le i\le4$, then $r_4$ belongs to $\mathcal C$,
to $\mathcal E$, or to the finite set
$\{(1,1,1,1)/3,(3,3,3,1)/5,(3,1,1,1)/5\}$.
The last two points form a two-cycle and the first is fixed.
\item[(ii)] Let $0<\eta\le10^{-4}$, $5/7-\eta\le a\le3/4$,
$0\le b-c\le\eta$, $|c-2a+1|\le\eta$ and
$100\eta\le d\le1-a+\eta$. Then $a_6=a_3$.
\item[(iii)] Under (ii)'s hypotheses, if also $a\ge5/7$ and $a_6=a$,
then $\nu(r_6)\ge\nu(r_3)$, where $\nu=b-c+|c-2a+1|$.
\item[(iv)] If $\eta>0$, $W(r)\le\eta$, $\min D(r)\ge100\eta$, and
$r_9\notin\mathcal C$, then $r_9-r_6=r_6-r_3$.
\end{enumerate}
Membership of each exceptional phase is a linear predicate, obtained
by eliminating its displayed parameters. For example, on the second
phase require $5/7\le a\le3/4$, $c=3a-2$, $d=3-4a$, and
$1-a\le b<2-2a$. On the third require the same range for $a$,
$b=a$, $c+d=2-2a$, and $0\le d<1-a$.
The encoding is \file{nine_global_entry_audit.py}. In the ancillary
directory run \runcheck{entry}. The four negations must be
unsatisfiable in both Z3 5.4.1 and cvc5 1.4.2. Direct substitution
verifies the cycles and the identity of the third iterate on the core;
\file{nine_entry_checks.py} also certifies the core inequalities.

\section{Finite core and lattice calculations}\label{app:routes}
The affine checks use \runcheck{routes}. The probability tables use
\runcheck{weights}. These calculations
use Python exact integers and rational numbers; SymPy 1.14.0 handles
symbolic arithmetic and Z3 5.4.1 supplies rational witnesses for
implications that are subsequently checked by exact identities.

\paragraph{Sorting words.}
Start with the ordered list in \eqref{eq:borrow}, indexed from zero.
For a permutation $\pi$, impose each adjacent inequality
$d_{\pi_i}-d_{\pi_{i+1}}\ge\boldsymbol1_{\pi_i>\pi_{i+1}}$.
Subtract symmetric sorted slots to form the new affine vector. Compose
both guards and output matrices along a word. If on a phase the result
is $q+\delta$, all guards at repetition $j$ have the form
$g(q,B)+jg_{\rm lin}\delta\ge0$. Over the integers the maximal number
of repetitions is the minimum of $1+\lfloor g(q,B)/(-g_{\rm lin}\delta)\rfloor$
over guards with negative slope. Restrict homogeneous guards to vertices
of the incoming face, enumerated by four independent equality rows of
\eqref{eq:core}. For an identically zero leading guard, substitute the
integer phase offsets; a positive leading margin absorbs every bounded
error on compact relative interiors. Equivalently a nonnegative rational
combination of the defining inequalities certifies a guard. Apply these
tests to every intervening step of the word.

\paragraph{Facet and corner calculations.}
Intersect each decreasing facet with the current face. Enumerate its vertices
as above, retain intersections of relative codimension one, and pull back
the first-exit comparisons $h/s<h'/s'$ through the preceding moves. This
gives Table~\ref{tab:routes}. The two releases out of a higher-codimension
intersection require the following extra statements. At $B_0=A_0$ the
four recurrent $B_0$ offsets each have eight successive stopping remainders.
Each reaches the $A_2$ return in at most four steps and that return increases
$B_2$ by four. At $B_0=B_1=A_1$ the twelve recurrent offsets each have
sixteen successive remainders; each reaches $AB$ within 23 steps and its
return increases $B_2$ by eighteen. These are the two points where
additional vanishing leading guards must be checked along the released
direction.

For clarity, the integer iteration producing the return words can be
reconstructed without a list of input strings. Use exact slack coordinates
$p=A_2$, $r=A_1$ at $AA$, $p=A_2$, $r=B_1$ at $AB$, and $p=B_0$,
$r=B_1$ at $BB$. Infer each three-step affine update by the sorting rules
above on the relative interior, retaining constant guards on $p,r$. On an
$A_2$ face, $p$ maps to $2-p$ if $p<0$, $2-2p$ if $0\le p<2$, and $p-4$
otherwise. Its recurrent states are $-2,-1,0,2,3,4$. On a $B_0$ face
$p'=|p+1|-4$, entering the cycle set $\{-4,-3,-2,-1\}$.
At $AB$ the $A_2$ update drives $r'=|r+2|-5+(-p)_+$ into $[-5,0]$.
At $BB$, $r'=|r+(-p-1)_++2|-5$. After $p$ is recurrent, the increment
on the upper branch is nonpositive over a full $p$ period with strictly
negative total, and the lower branch folds upward to at least $-5$.
At $AA$ both coordinates after a return are at least $-2$ and
$\max(p,r)\le M$, $M\ge5$, implies that their maximum two returns later
is at most $M-1$ on $p,r\ge-2$. These last inequalities are verified
on all pairs of three-step sorting cells by nonnegative linear combinations
of their guards. Iterating the resulting finite sets gives the return
lengths and displacements in the table.
The encodings for these assertions are \file{nine_entry_checks.py},
\file{nine_a_facet_checks.py}, \file{nine_b_facet_checks.py},
\file{nine_aa_checks.py}, \file{nine_ab_checks.py}, \file{nine_bb_checks.py},
\file{nine_b_routes_checks.py}, and \file{nine_releases_checks.py}.

\paragraph{End of each return.}
For each recurrent phase compute the first failed guard as above. At
$AA$ the outgoing $A_0$ remainders occupy intervals of length $30$ or
$5$ on the two phase cycles, giving $435$ possible terminal inputs.
In the chart $q=BQ(U/B)+(0,i,j,l)$ the terminal update is
\[
 (i',j',l')=(2t+(-l-1)_+-(l-1)_+,\,3t+\min(0,-l-1)+i+1,\,i-j-t).
\]
Here $t=\min(l,1)$ and $U'=U-t$. Iterate the reached triples to their first repeats.
The periods are $1,2,6,32,37$, with respective longitudinal shifts $0,0,0,-2,-3$. At $AAB$ there are sixteen input phases with $24$ possible remainders
each. The chart is
\[
 q=(U,2U-B-p,3U-2B-p-e,(3U-2B-p-s)/2)
\]
with $p=A_2$, $e=A_1$, $s=B_1$ and $U\equiv p+s\pmod2$.
Follow the digits until either a recurrent $AA$ word starts or a
nine-step translation starts on
\[
 (U,2U-B+1,3U-2B+3,(3U-2B+1-s)/2),
 \qquad s\in\{-2,-1,0,1\}.
\]
It decreases $U$ by two. On the moving $AAA$ and $AAB$ returns check
all word inequalities on $4B/5+252\le U\le5B/6-256$, $B\ge10^6$.
Write $B=5b+\beta$ and stop at the unique $U=4b+e$ with
$256\le e<256+|\Delta U|$ in the reached congruence class.
Every remaining step is affine in $b$ with fixed bounded offset.
Follow the offsets to their first repeat and check their sorting
inequalities for every integer $b\ge200000$.

At $BBB$ write
$q=(3V-B+i,V,B-V+j,2V-B+l)$ for $3B/5<V<2B/3$.
The first coordinate falls by one per step. In the coordinates
$(x,y_0,z_0)=(j+1,i-2l+3,2i+2j-l+4)$ the offset map is
\[
 (x,y_0,z_0)\longmapsto(y_0-\min(x,0)-1,z_0-1,|x|-1).
\]
The sum is $2$ modulo three. Three steps suffice for the lower bounds
$(-3,-2,-1)$, which are invariant. In that orthant, for $M\ge9$ a
maximum at most $M$ becomes at most $M-1$ after six steps, while a
maximum at most eight remains at most eight after six. Verify these
inequalities on all six-sign words by nonnegative combinations of affine
guards. Iterate the resulting finite integer box through intervening
steps to a nine-cycle of offsets. The representative
$(i,j,l)=(-3,2,0)$ returns with $V$ decreased by three, under
$3B/5+256\le V\le2B/3-256$ and $B\ge10^6$. Stop at
$V=3b+e$, $e\in\{256,257,258\}$, and verify the fifteen remaining
affine trajectories to first repetition.

To certify first capture, compare every intermediate affine state of
the final incoming word with each member of the terminal cycle, allowing
an arbitrary number of earlier translations. The first coordinate fixes
the only possible repetition number; the other coordinates decide equality.
On stationary $AAA$ cycles the offsets $(i,j,l)$ equal to
$(2,4,-1)$, $(0,-2,-1)$, $(-2,0,-1)$, $(-2,-4,-1)$ move capture five steps
earlier, and $(2,2,1)$, $(0,-2,1)$, $(-2,-2,3)$, $(-2,-4,3)$ move it two
steps earlier. All other stationary arrivals occur as counted. For a
drifting arrival, the minimum prior first coordinate exceeds the maximum
cyclic first coordinate on compact strict source cells at sufficiently
large $B$; minimize an affine function of the repetition count at either
end. This also excludes capture in earlier route segments. The terminal
encodings are \file{nine_aaa_checks.py}, \file{nine_aab_checks.py} and
\file{nine_bbb_checks.py}, with the word reconstruction in
\file{nine_phase_checks.py}. The full route composition and all first-exit
competitors are checked by \file{nine_core_checks.py}.

\paragraph{Weights and covariance.}
For each route and each $\beta$, begin with the full lattice.
At each stopping interval solve the single linear Diophantine equation
in Section~\ref{sec:law} by column reduction, retain every soluble
remainder and apply its exact conditional weight. Propagate the affine
clock and terminal maps. Add weights only for coincident observables
after retaining the full image lattice at every intermediate step.
The checker \file{nine_kernel_checks.py} reconstructs every entry of
\file{nine_kernel_results.json} and checks that the total on every
route is one. The kernels have affine clocks and terminal coordinates up
to bounded errors on compact subsets. This follows from the bounded
remainders, finitely many returns per route and the exact terminal words.
The paired-level hypothesis on each core sorting word is checked by
\file{odd_base_covariance_checks.py} in \runcheck{covariance}.

\section{Zero returns and the nine-step descent}\label{app:clock}
All variables in this section are integers. Write $q=(u,v',w,z)$.
For an input with only $r$ positive differences, use the borrowed
formula for the first $r$ coordinates with $9-2r$ copies of the
central digit $B-1$. Sort and take four outer differences; at the zero
vector the update is zero. Thus every step can be encoded using integer
linear arithmetic, comparisons and conditional expressions.

\paragraph{First access.}
Zero at the next step means that the three central sorted digits agree.
Group the digits of a positive source into $B-1$ and the pairs of sums
$B,B-1,B-1,B-2$ from \eqref{eq:borrow}. The common digit must lie
between $(B-2)/2$ and $B/2$, by counting the entries on either side.
For $B=2a+1$, counting above and below $a$ gives at least two coordinates
equal to $a$ among $v',w,z$ as the necessary and sufficient condition.
For $a\ge16$ the negation of the following implication is unsatisfiable:
after three updates through positive states from a positive source, this
gate is false. The encoding is \file{odd_zero_access_audit.py}
(Z3 5.4.1 and cvc5 1.4.2).
For $B=2a$, the common value must be $a$; at $a-1$ the outer and two
middle pairs and central digit already supply four larger entries.
Counting at $a$ gives the gate $u=a, v'\ge a-1,z<a$.
The first output coordinate equals $\max(u-1+\boldsymbol1_{u=v'},B-z)$.
For a positive predecessor with output first coordinate $a$ the source
can therefore only be $(a,a,a,a)$ or $(a+1,a,a,a)$, whose outputs
have second coordinate one or two. For $B\ge40$ neither is in the gate.

\paragraph{Cover of returning zero states.}
Require $B\ge40$, $0<u<B$, $u\ge v'\ge w\ge z=0$. Rank one is included
in full. At rank two ($v'>0=w$) the cover is the union of
\[
 |u-v'|\le3,\quad |u+v'-B|\le3,\quad |2u-B|\le4,\quad
 |2v'-B|\le4,\quad B-u\le3 .
\]
At rank three take the union of the following five rows. A comma inside
the right column separates alternatives.
\begin{center}\small
\begin{tabular}{@{}ll@{}}
\toprule
Required & Additional condition\\\midrule
$v'=w$ & $|u-v'|\le2$, $|u+v'-B|\le2$, $|2v'-B|\le3$\\
$u=v'$ & $|u+w-B|\le2$, $|2u-B|\le3$\\
$|2u-B|\le3$ & $|2v'-B|\le3$\\
$B-u\le3$ & $|2v'-B|\le3$, $v'=w$, $|v'+w-B|\le3$\\
$|4u-3B|\le8$ & $|4y-jB|\le12$ for some $y\in\{v',w\}$, $1\le j\le3$\\
\bottomrule
\end{tabular}
\end{center}
At a fixed base the rank bounds are $B,24B$ and
$(12+9+16+36+210)B$, by fixing an unrestricted coordinate and counting
integers in each restricted interval. For both parities and ranks two
and three impose a first return at time $j=1,2,3,4$ with all intermediate
inner differences positive. Add the negation of the indicated union.
Every such query is unsatisfiable. In the same encoding require $q_4>0$,
$B\ge2$ and test the four negations of $S_i(g_B(F_Bq))\le S_i(g_B(q))+1$.
The file \file{nine_zero_return_audit.py} runs these checks with
Z3 5.4.1 and cvc5 1.4.2, including a direct test of the $+8$ dispersion
bound. Run \runcheck{zeros} for the access, cover and dispersion
calculations.

\paragraph{Nine-step descent.}
Fix one of the $5!$ weak ascending orders of the source gaps, and impose
one of $B-M_g-4m_g>100$, $B-6m_g>100$. These alternatives cover $P>100$.
At a node of depth $i$ retain the ancestral guards, positivity,
and $G(q_0)\le G(q_i)\le G(q_0)+8i$. A child supplies descending orders
for the eight digits other than the largest $B-1$ and for the five output
gaps. Using the integer tie rule of Appendix~\ref{app:routes} gives the
child's affine state, guards and dispersion. For freshly sorted
eight digits $d_0\ge\cdots\ge d_7$ the successor is
$(B-1-d_7,d_0-d_6,d_1-d_5,d_2-d_4)$. For each node ask whether a
positive successor with $G\ge G(q_0)$ can avoid all listed child cells;
at a leaf ask whether any such successor exists. Sorting in this query
uses min/max comparators independently of the child orders.
The upper bound used in pruning follows from the one-step $+8$ inequality.

The trees in \file{nine_clock_certificates.json} have 29,534 coverage nodes
and depth at most eight. The verifier rebuilds the states from parent
indices and sorting permutations, checks connectivity and the full inventory
of source cases, and requires unsatisfiability of every coverage query
in Z3 5.4.1 and Yices 2.7.0. Induction down a tree excludes a nine-step
counterexample to \eqref{eq:descent}. Run \runcheck{descent}
(\file{nine_clock_audit.py} and \file{nine_clock_checks.py}).

\section*{Verification}
The supplementary programs reconstruct the rational and integer calculations
from the digit rule. Solver replays use two decision procedures on each
universal query. A separate direct column-subtraction test on one inverse
branch of the distinguishing route uses five fresh consecutive bases and
compares the original-number time and terminal differences with all four
predicted atoms. Each observed trajectory matched a unique atom. The largest
absolute frequency discrepancy over the five bases was below $0.012$; the
largest normalized clock error was below $0.00014$. The finite tests concern
the indicated branch; the argument
for arbitrary bases is the lattice and entry proof above.

An independent re-derivation of every reduction from the text of an earlier
draft of this paper was also carried out, with a separately written transition
model (column subtraction and a sorting network) and separate Z3 encodings.
Z3 found no counterexample to the four real entry implications, the core and
three-cycle identities, or the stated one-step bounds, first-access and sixteen
zero-return implications; the integer queries imposed no upper bound on $B$.
Direct subtraction at bases not used here tested the separating route of
Section~4.4 and the 165 kernels on about 2.6 million conditional draws, and
the differences between the five laws on 1.6 million uniform strings; the
pooled comparisons and targeted replications were consistent with the
predictions. These tests are statistical. The nine-step descent \eqref{eq:descent} held at every
qualifying state for $101\le B\le160$ and at five larger bases, and on
3.1 million sampled states with $B$ up to $2\cdot10^9$. These encodings and
their logs are part of the ancillary material.

\section*{Data availability}
The ancillary material is archived at Zenodo as version~1.3 of \href{https://doi.org/10.5281/zenodo.24325109}{doi:10.5281/zenodo.24325109}, and maintained at \url{https://github.com/galyavin/kaprekar-dynamics}.

\section*{Tools}
The research, computations and drafting were carried out with substantial
assistance from AI systems (OpenAI Codex and Anthropic Claude), under the
author's direction and close supervision. The computations use Python, SymPy,
Z3, cvc5, Yices and C++. The author takes responsibility for the contents.

\begingroup\raggedright
\begin{thebibliography}{Thak}

\bibitem[AbSZT]{major}
R.~Abdulghafor, F.~Shahidi, A.~Zeki and S.~Turaev, \emph{Dynamics of doubly stochastic quadratic operators on a finite-dimensional simplex}, Open Math. \textbf{14} (2016), 509--519; \href{https://doi.org/10.1515/math-2016-0045}{doi:10.1515/math-2016-0045}.

\bibitem[Dah]{Dahl}
C.~D. Dahl, \emph{Coarse-grained drift fields and attractor-basin entropy in Kaprekar's routine}, Entropy \textbf{28} (2026), 92; \href{https://doi.org/10.3390/e28010092}{doi:10.3390/e28010092}.

\bibitem[G1]{G1}
R.~Galyavin, \emph{Cycles of the four-digit Kaprekar map in even bases}, preprint (2026), \href{https://doi.org/10.5281/zenodo.24325110}{doi:10.5281/zenodo.24325110}.

\bibitem[G2]{G2}
R.~Galyavin, \emph{Arrival times for the four-digit Kaprekar map in even bases}, preprint (2026), \href{https://doi.org/10.5281/zenodo.24329321}{doi:10.5281/zenodo.24329321}.

\bibitem[G3]{G3}
R.~Galyavin, \emph{A common arrival-time law for five-digit Kaprekar dynamics}, preprint (2026), \href{https://doi.org/10.5281/zenodo.24330159}{doi:10.5281/zenodo.24330159}.

\bibitem[KD]{kay}
A.~Kay and K.~Downes-Ward, \emph{Fixed points and cycles of the Kaprekar transformation: 1. Odd bases}, J. Integer Seq. \textbf{25} (2022), Article 22.6.7.

\bibitem[Thak]{thakur}
D.~S. Thakur, \emph{Kaprekar phenomena}, in: Topics in Number Theory, RMS Lecture Notes Series \textbf{26}, Ramanujan Mathematical Society, 2020, 61--70.

\end{thebibliography}
\endgroup
\end{document}
